\documentclass{article}

\usepackage{amsmath,amssymb}
\usepackage{xurl}
\usepackage[hidelinks]{hyperref}
\setlength{\emergencystretch}{2em}

% Shared commands for notes in docs/paper-gaps/.

\DeclareUnicodeCharacter{2080}{\ifmmode{}_{0}\else\textsubscript{0}\fi}
\DeclareUnicodeCharacter{2081}{\ifmmode{}_{1}\else\textsubscript{1}\fi}
\DeclareUnicodeCharacter{2082}{\ifmmode{}_{2}\else\textsubscript{2}\fi}
\DeclareUnicodeCharacter{2099}{\ifmmode{}_{n}\else\textsubscript{n}\fi}

\newcommand{\Lean}{\textsc{Lean}}
\ExplSyntaxOn
\NewDocumentCommand{\leanid}{m}{
  \ifmmode
    \textsf{\detokenize{#1}}
  \else
    \group_begin:
    \urlstyle{sf}
    \tl_set:Nn \l_tmpa_tl {#1}
    \tl_replace_all:Nnn \l_tmpa_tl {\_} {_}
    \exp_args:NV \path \l_tmpa_tl
    \group_end:
  \fi
}
\ExplSyntaxOff
\providecommand{\tr}{\operatorname{tr}}
\newcommand{\tnleanIssueBase}{https://github.com/LionSR/TNLean/issues/}
\newcommand{\tnleanPrBase}{https://github.com/LionSR/TNLean/pull/}
\newcommand{\tnleanDocBase}{https://sirui-lu.com/TNLean/docs/}
\newcommand{\ghissue}[1]{\href{\tnleanIssueBase#1}{GitHub issue \##1}}
\newcommand{\ghpr}[1]{\href{\tnleanPrBase#1}{PR \##1}}
\newcommand{\doclink}[1]{\href{\tnleanDocBase#1.html}{\path{#1}}}

% Dirac notation and the identity operator, as in blueprint/src/macros/common.tex.
% \providecommand keeps note-local definitions authoritative.
\usepackage{dsfont}
\providecommand{\ket}[1]{|#1\rangle}
\providecommand{\bra}[1]{\langle #1|}
\providecommand{\braket}[2]{\langle #1 | #2 \rangle}
\providecommand{\ketbra}[2]{|#1\rangle\!\langle #2|}
\providecommand{\Id}{\mathds{1}}
\providecommand{\one}{\mathds{1}}
\providecommand{\C}{\mathbb{C}}
\providecommand{\MN}[1]{M_{#1}(\C)}
\providecommand{\MD}{M_D(\C)}

% External referencing for paper sources.  The build helper writes sanitized
% auxiliary files under build/paper-aux/ and excludes duplicated source labels.
\usepackage{xr}

% Import a generated paper auxiliary file with a caller-supplied label prefix.
% The candidate paths support builds from docs/paper-gaps/, the repository root,
% and build/paper-aux/ itself.
\newcommand{\importpaperaux}[2]{%
  \IfFileExists{../../build/paper-aux/#2.aux}{%
    \externaldocument[#1]{../../build/paper-aux/#2}%
  }{%
    \IfFileExists{build/paper-aux/#2.aux}{%
      \externaldocument[#1]{build/paper-aux/#2}%
    }{%
      \IfFileExists{#2.aux}{%
        \externaldocument[#1]{#2}%
      }{}%
    }%
  }%
}

% General external reference macros
\newcommand{\paperref}[2]{\ref{#1-#2}}
\newcommand{\papereqref}[2]{(\ref{#1-#2})}

% CPSV16 source (1606.00608 / MPDO-22-12-17-2.tex) external document and shortcuts
\importpaperaux{cpsv16-}{1606.00608/MPDO-22-12-17-2-tnlean}
\newcommand{\cpsvref}[1]{\paperref{cpsv16}{#1}}
\newcommand{\cpsveqref}[1]{\papereqref{cpsv16}{#1}}

% Verdict marker: \gapnote{<kind>}{<status>}. Kind is the policy
% classification (clarification, local-correction, scope-restriction,
% unfaithful, false-source, open-gap); status is open, wip (elimination
% underway), resolved, or historical. Machine-read by the site index;
% prints a verdict line.
\newcommand{\gapnote}[2]{\par\noindent\textbf{Verdict:} #1 (#2).\par}


\title{Vertical Canonical-Form Grouping in Proposition 4.13}
\author{}
\date{2026-07-27}

\begin{document}
\maketitle
\gapnote{clarification}{resolved}

\section{Source statement and proof order}

Proposition~4.13, whose source label is \path{Prop:IV.12}, of
Ref.~\cite{Cirac2016MPDO_arXiv} states that a canonical-form tensor generating
matrix product density operators is also in canonical form in the vertical
direction. Let $\widetilde M$ denote the tensor viewed vertically: its letters
are indexed by pairs of horizontal bond indices, and its matrices act on the
physical space. With the rectangular map called an isometry in the source, the
claimed decomposition is
\begin{align}
    U\widetilde M U^\dagger
        &=\bigoplus_\alpha\mu_\alpha\otimes M_\alpha.
    \label{eq:cpgsv17_vertical_cf_grouping_source_vertical_form}
\end{align}
Here each $\mu_\alpha$ is positive and diagonal, and the tensors $M_\alpha$
form a basis of normal tensors. In this rectangular orientation, $U$ maps the
physical space onto the retained sector space and is therefore a coisometry:
\begin{align}
    UU^\dagger
        &=\Id.
    \label{eq:cpgsv17_vertical_cf_grouping_source_coisometry}
\end{align}
The exact vertical canonical form also includes the reverse reconstruction
\begin{align}
    \widetilde M
        &=U^\dagger
        \left(\bigoplus_\alpha\mu_\alpha\otimes M_\alpha\right)U.
    \label{eq:cpgsv17_vertical_cf_grouping_source_vertical_reconstruction}
\end{align}
Proposition~4.13 calls $U$ an isometry and prints only
Eq.~\ref{eq:cpgsv17_vertical_cf_grouping_source_vertical_form}. The
orientation in
Eq.~\ref{eq:cpgsv17_vertical_cf_grouping_source_coisometry} and the reverse
identity in
Eq.~\ref{eq:cpgsv17_vertical_cf_grouping_source_vertical_reconstruction} give
the source-faithful rectangular formulation recorded in
\path{cpgsv17_vertical_tensor_definition.tex}.

The proof of Proposition~4.13 has a fixed order. It first combines positivity
of the generated operators with Lemma~L of
Ref.~\cite{Cirac2016MPDO_arXiv} to exclude one-sided transitions and
nontrivial periodic sectors. Only after obtaining vertical canonical form does
it write the vertical decomposition as
\begin{align}
    \bigoplus_{\alpha,k}
        \mu_{\alpha,k}X_{\alpha,k}M_\alpha X_{\alpha,k}^{-1}.
    \label{eq:cpgsv17_vertical_cf_grouping_source_sector_form}
\end{align}
The rest of the source proof shows that $\mu_{\alpha,k}>0$ and that every
$X_{\alpha,k}$ is proportional to a unitary. The copy label $k$ in
Eq.~\ref{eq:cpgsv17_vertical_cf_grouping_source_sector_form} then becomes the
diagonal index of $\mu_\alpha$ in
Eq.~\ref{eq:cpgsv17_vertical_cf_grouping_source_vertical_form}.

Thus the positive diagonal coefficients arise from the vertical canonical
decomposition. They do not arise by flattening scalar weights from a
previously selected horizontal decomposition while preserving diagonal
restrictions of its blocks. The relevant source passage is
\path{Papers/1606.00608/MPDO-22-12-17-2.tex:1861--1922}.

\section{Exact finite-sum statement}

Let $(B_k)_{k\in I}$ be a finite indexed family of blocks already grouped into
copies of representatives $(A_\alpha)_{\alpha\in J}$. For each $\alpha$, let
$r_\alpha$ be the number of copies, equivalently the size of the diagonal
matrix $\mu_\alpha$. The trace coefficient
$m_\alpha=\tr(\mu_\alpha)$ is a different quantity, as explained in
\path{cpgsv17_vertical_tensor_definition.tex}. Let
\begin{align}
    e:I
        &\simeq
        \{(\alpha,j):\alpha\in J,\ 0\leq j<r_\alpha\}
    \label{eq:cpgsv17_vertical_cf_grouping_grouping_equivalence}
\end{align}
be the grouping bijection. Suppose that, for every $k\in I$,
\begin{align}
    \dim B_k
        &=\dim A_{e(k)_1},
    \label{eq:cpgsv17_vertical_cf_grouping_pointwise_dimension}\\
    \nu_k
        &=\omega_{e(k)},
    \label{eq:cpgsv17_vertical_cf_grouping_pointwise_weight}\\
    V^{(N)}(B_k)
        &=V^{(N)}(A_{e(k)_1})
        \qquad (N\geq1).
    \label{eq:cpgsv17_vertical_cf_grouping_pointwise_mpv}
\end{align}
Finite-sum reindexing then gives, for every $N$, including $N=0$,
\begin{align}
    V^{(N)}\!\left(\bigoplus_{k\in I}\nu_kB_k\right)
        &=
        V^{(N)}\!\left(\bigoplus_{\alpha\in J}
        \bigoplus_{j=0}^{r_\alpha-1}
        \omega_{\alpha,j}A_\alpha\right).
    \label{eq:cpgsv17_vertical_cf_grouping_regrouped_mpv}
\end{align}
At $N=0$, this equality follows from
Eq.~\ref{eq:cpgsv17_vertical_cf_grouping_pointwise_dimension}. Summing those
dimension equalities along the bijection in
Eq.~\ref{eq:cpgsv17_vertical_cf_grouping_grouping_equivalence} gives equality
of the total bond dimensions.

The formal theorem is
\path{MPOTensor.sameMPV2_toTensorFromBlocks_verticalAssembledTensor_of_equiv}.
In this rendered identifier, the ordinary digit~$2$ denotes the subscripted
character used in the source declaration.

\section{Boundary of the result}

The normalized BNT-refined horizontal predicate
\path{MPOTensor.IsHorizontalCF} is strictly stronger than the literal CPSV
canonical form used in Proposition~4.13 of
Ref.~\cite{Cirac2016MPDO_arXiv}. Its retained sectors fill the horizontal bond
space, every repeated-copy weight is nonzero, and its normal blocks are already
grouped by BNT representative. The theorem under this stronger hypothesis is
preserved, while the literal implication has now been proved independently.
Under literal CPSV canonical form, the completed proof includes periodic
exclusion, normal vertical block decomposition, finite-chain separation,
vertical phase-class grouping with positive effective coefficients, the
actual-grouped Figure~8 identity, unitary gauge normalization, normalized
physical-sector maps, and the final rectangular coisometry.

The periodic-exclusion component is recorded by
\path{MPSTensor.IsCPSVCanonicalForm.exists_not_commute_of_displaced},
\path{MPOTensor.hasNoPeriodicVectors_verticalTensor_of_cpsvCanonicalForm}, and
the stronger-hypothesis variant
\path{MPOTensor.hasNoPeriodicVectors_verticalTensor_of_horizontalCF}. Its
quantifier correction is documented in
\path{docs/paper-gaps/cpgsv17_periodic_sector_projector.tex}. The subsequent
representative-sector decomposition is
\path{MPSTensor.CPSVCanonicalFormData.BNTRefinement.groupedTensor_sameMPV₂Pos_representativeSectorDecomposition}.

An earlier auxiliary route compared one scalar $\mu_\alpha$ with several
weights $\omega_{\alpha,j}$ under the assumptions
\begin{align}
    \mu_\alpha^N
        &=\sum_j\omega_{\alpha,j}^N
        \qquad (N\geq1).
    \label{eq:cpgsv17_vertical_cf_grouping_power_sum_assumption}
\end{align}
Equation~\ref{eq:cpgsv17_vertical_cf_grouping_power_sum_assumption} gives a
valid conditional comparison at positive lengths, but it is not the flattening
argument in Proposition~4.13. The corresponding terminal lemmas have therefore
been removed. In the source proof, the several weights already belong to the
separate vertical sectors in
Eq.~\ref{eq:cpgsv17_vertical_cf_grouping_source_sector_form}; no single
horizontal coefficient is replaced by their power sum.

Diagonal restriction also need not preserve injectivity of a horizontal block.
For $d=D\geq2$, let $E_{ij}$ be the matrix units and define
\begin{align}
    B^{(i,j)}
        &=d^{-1/2}E_{ij}.
    \label{eq:cpgsv17_vertical_cf_grouping_matrix_unit_block}
\end{align}
These matrices satisfy
\begin{align}
    \operatorname{span}
        \{B^{(i,j)}:1\leq i,j\leq d\}
        &=M_D(\C),
    \label{eq:cpgsv17_vertical_cf_grouping_full_matrix_span}\\
    \sum_{i,j}(B^{(i,j)})^\dagger B^{(i,j)}
        &=\Id.
    \label{eq:cpgsv17_vertical_cf_grouping_matrix_unit_normalization}
\end{align}
Their diagonal restriction is
\begin{align}
    B^{(i,i)}
        &=d^{-1/2}E_{ii},
    \label{eq:cpgsv17_vertical_cf_grouping_diagonal_restriction}
\end{align}
whose span is only the diagonal subalgebra. Therefore, the BNT grouping in the
conclusion must come from the new vertical canonical decomposition; one cannot
obtain it by declaring the diagonal restrictions of the horizontal blocks to
be normal blocks.

The exact reindexing identity
Eq.~\ref{eq:cpgsv17_vertical_cf_grouping_regrouped_mpv} is complete once the
bijection and pointwise block identifications are available. Under
\path{MPOTensor.IsHorizontalCF}, the one-sided invariant-matrix calculation at
source lines~1874--1887 proves invariant-projector closure for the vertical
tensor:
\path{MPOTensor.IsHorizontalCF.hasInvariantProjectorClosure_verticalTensor}.
Under literal CPSV canonical form, the corresponding result is
\path{MPSTensor.IsCPSVCanonicalForm.hasInvariantProjectorClosure_verticalTensor}.

Under the closure hypothesis, the general decomposition theorem
\path{MPOTensor.IsHorizontalCF.exists_irreducible_verticalBlockDecomp_with_isometry}
produces a complete orthogonal family of irreducible reducing corners,
isometries, both intertwining identities, exact compression, and equality of
matrix product vectors at every length. Its range projections are complete,
pairwise orthogonal, and commute with every vertical letter, as recorded by
\path{MPOTensor.IsHorizontalCF.exists_complete_verticalReducingProjectors}.

The nonzero corners can be spectrally normalized without changing their
physical embeddings. The resulting blocks are normal, their weights are
positive, both intertwining identities and exact compression remain valid, and
the retained sectors reconstruct every vertical letter exactly. This is
\path{MPOTensor.IsHorizontalCF.exists_normal_verticalBlockDecomp_with_isometry}
in \path{TNLean/MPS/MPDO/VerticalSpectral.lean}. The orthogonal complement of
the retained range projections consists exactly of the omitted zero corners;
removing those corners therefore introduces no false completeness assertion.

The same normal-corner result is complete under literal CPSV canonical form.
The theorem
\path{MPSTensor.IsCPSVCanonicalForm.exists_normal_verticalBlockDecomp_with_isometry}
in \path{TNLean/MPS/MPDO/CPSVPeriodicExclusion.lean} combines literal
invariant-projector closure with literal periodic exclusion. It produces
positive-dimensional normal blocks $B_k$, positive weights $\nu_k$, and
pairwise orthogonal isometries $V_k$ satisfying, for every vertical letter
$v$,
\begin{align}
    \widetilde M_vV_k
        &=V_k(\nu_kB_k^v),
    \label{eq:cpgsv17_vertical_cf_grouping_literal_right_intertwining}\\
    V_k^\dagger\widetilde M_v
        &=(\nu_kB_k^v)V_k^\dagger,
    \label{eq:cpgsv17_vertical_cf_grouping_literal_left_intertwining}\\
    \nu_kB_k^v
        &=V_k^\dagger\widetilde M_vV_k,
    \label{eq:cpgsv17_vertical_cf_grouping_literal_compression}\\
    \widetilde M_v
        &=\sum_kV_k(\nu_kB_k^v)V_k^\dagger.
    \label{eq:cpgsv17_vertical_cf_grouping_literal_reconstruction}
\end{align}
The omitted orthogonal complement is an all-zero sector. This theorem does not
itself partition the $B_k$ into gauge-equivalent representative families; the
subsequent grouping, normalization, and coisometry theorems complete
Proposition~4.13.

The literal sector-compression separation theorem
\path{MPSTensor.IsCPSVCanonicalForm.exists_sectorCompression_ne_zero_of_corner}
is proved in \path{TNLean/MPS/MPDO/CPSVVerticalBNT.lean}. For every matrix
$P$, it states
\begin{align}
    \left(\exists v:\ P\widetilde M_vP\ne0\right)
        &\Longrightarrow
        \left(\exists N\geq0:\
        P_1H^{(N+1)}P_1\ne0\right).
    \label{eq:cpgsv17_vertical_cf_grouping_sector_compression_separation}
\end{align}
If every compression on the right vanished, the two-sided first-site insertion
by $P$ would agree with zero on every positive-length matrix product vector.
Original-space Lemma~L would then force the inserted tensor, and hence every
corner $P\widetilde M_vP$, to vanish. This implication uses neither positivity
nor vertical phase-class grouping. It proves only the separation clause at
source line~1898. The subsequent trace argument cannot be applied until the
source sectors have actually been grouped.

The normalized corners are grouped by matrix-product-vector phase class under
literal CPSV canonical form in
\path{MPSTensor.IsCPSVCanonicalForm.exists_verticalBNTGrouping_with_isometry},
proved in \path{TNLean/MPS/MPDO/CPSVVerticalBNT.lean}. The wrapper under the
normalized BNT-refined horizontal hypothesis remains
\path{MPOTensor.IsHorizontalCF.exists_verticalBNTGrouping_with_isometry}.
The construction chooses one normal representative from each phase class.
Every copy has the representative's bond dimension and is a unit-modulus
scalar times an invertible conjugate of that representative. The scalar for
the chosen representative is one. The representatives are pairwise
gauge-phase distinct and form a basis of normal tensors for the weighted block
tensor.

The physical isometries remain unchanged by this grouping. Both intertwining
identities, exact compression, and literal reconstruction of each vertical
letter therefore survive. Each grouped sector has a physical range projector
that satisfies the representative-loop identity and gives a nonzero
compression at some finite length. The MPDO trace argument then proves that
every grouped coefficient is positive. This implements the phase-class
decomposition and positivity argument at source lines~1898--1902 by combining
the BNT characterization at source lines~1135--1148 with the preceding
isometry-preserving vertical-sector theorem. Gauge normalization and the
coisometry construction then give the positive-diagonal form asserted at
source lines~1895--1921. Equality of bond dimensions inside each phase class
follows from non-decay of the positive-length rectangular overlap, as in
Lemma~\texttt{equalMPS} of Ref.~\cite{Cirac2016MPDO_arXiv}; the comparison
does not use the empty-word coefficient.

Under either the literal CPSV hypothesis or the normalized BNT-refined
horizontal hypothesis, write $\nu_{\alpha,k}$ for a raw sector weight,
$\zeta_{\alpha,k}$ for its phase relative to the chosen representative, and
define
\begin{align}
    c_{\alpha,k}
        &=\nu_{\alpha,k}\zeta_{\alpha,k}.
    \label{eq:cpgsv17_vertical_cf_grouping_grouped_coefficient}
\end{align}
Let $P_{\alpha,k}$ be the physical range projector for the sector, and define
the representative-loop matrix by
\begin{align}
    (E_\alpha)_{ab}
        &=\tr(M_\alpha^{(a,b)}).
    \label{eq:cpgsv17_vertical_cf_grouping_representative_loop}
\end{align}
Inserting $P_{\alpha,k}$ into the one-site vertical loop selects
$c_{\alpha,k}E_\alpha$. This expression is independent of the copy index
because cyclicity of trace removes the internal similarity. Sector-compression
separation provides a nonzero compression, and MPDO positivity makes its trace
positive. With
\begin{align}
    \tr(c_{\alpha,k}E_\alpha)
        &=c_{\alpha,k}d_\alpha,
    \label{eq:cpgsv17_vertical_cf_grouping_sector_trace}\\
    c_{\alpha,1}
        &=\nu_{\alpha,1}>0,
    \label{eq:cpgsv17_vertical_cf_grouping_representative_normalization}
\end{align}
one obtains
\begin{align}
    d_\alpha
        &>0,
    \label{eq:cpgsv17_vertical_cf_grouping_positive_loop_trace}\\
    c_{\alpha,k}
        &>0.
    \label{eq:cpgsv17_vertical_cf_grouping_positive_grouped_coefficient}
\end{align}
Thus the coefficient denoted $\mu_{\alpha,k}$ in the source is
\begin{align}
    \mu_{\alpha,k}
        &=c_{\alpha,k}.
    \label{eq:cpgsv17_vertical_cf_grouping_source_grouped_coefficient}
\end{align}
These results are formalized in \path{TNLean/MPS/MPDO/VerticalBNT.lean},
\path{TNLean/MPS/MPDO/CPSVVerticalBNT.lean}, and
\path{TNLean/MPS/MPDO/SectorTrace.lean}.

The pairwise Gram-conjugation identity corresponding to the second displayed
diagram at source lines~1914--1919 is proved in
\path{TNLean/MPS/MPDO/FigureEightPairwise.lean}. Its dependent identification
with the actual grouped sector and distinguished copy is proved under the
normalized BNT-refined hypothesis in
\path{TNLean/MPS/MPDO/GroupedFigure8.lean}, and under literal CPSV canonical
form in \path{TNLean/MPS/MPDO/CPSVGroupedFigureEight.lean}. Neither proof
identifies either corner separately with the raw reflected adjoint.

The scalar-commutant theorem
\leanid{Kraus.IsNormal.eq_smul_one_of_commute} for the normal tensor
$M_\alpha$ makes the relative Gram ratio scalar. Positivity of the two Gram
matrices gives
\begin{align}
    X_{\alpha,k}^\dagger X_{\alpha,k}
        &=\omega_{\alpha,k}
        X_{\alpha,1}^\dagger X_{\alpha,1},
    \label{eq:cpgsv17_vertical_cf_grouping_gram_ratio}\\
    \omega_{\alpha,k}
        &>0.
    \label{eq:cpgsv17_vertical_cf_grouping_positive_gram_scalar}
\end{align}
Consequently, the relative gauge
\begin{align}
    \omega_{\alpha,k}^{-1/2}
        X_{\alpha,k}X_{\alpha,1}^{-1}
        &\quad\text{is unitary}.
    \label{eq:cpgsv17_vertical_cf_grouping_normalized_relative_gauge}
\end{align}
The grouping theorem exposes the normalization
\begin{align}
    X_{\alpha,1}
        &=\Id.
    \label{eq:cpgsv17_vertical_cf_grouping_distinguished_gauge}
\end{align}
Under the normalized BNT-refined hypothesis, the formal results are
\path{MPOTensor.IsMPDO.grouped_sector_gram_eq_pos_smul_one} and
\path{MPOTensor.IsMPDO.grouped_sector_exists_unitary_normalization} in
\path{TNLean/MPS/MPDO/GroupedGramNormalization.lean}. Their literal-CPSV
counterparts are
\path{MPSTensor.IsCPSVCanonicalForm.grouped_sector_gram_eq_pos_smul_one} and
\path{MPSTensor.IsCPSVCanonicalForm.grouped_sector_exists_unitary_normalization}
in \path{TNLean/MPS/MPDO/CPSVGroupedGramNormalization.lean}.

The reflected marked-chain argument, two-sector separation, and dependent
dimension identification at source lines~1909--1919 are therefore complete.
So are the positive scalar Gram identity and the unitary normalization of
every grouped gauge. Under the normalized BNT-refined hypothesis, the
normalized gauges are composed with the physical sector isometries in
\path{TNLean/MPS/MPDO/NormalizedGroupedSectors.lean}. The literal theorem
\path{MPSTensor.IsCPSVCanonicalForm.exists_normalized_grouped_sector_maps} is
proved in \path{TNLean/MPS/MPDO/CPSVNormalizedGroupedSectors.lean}. In both
cases, the maps $W_{\alpha,q}$ have mutually orthogonal isometric ranges,
satisfy the representative intertwining identities, and reconstruct every
vertical letter exactly.

The representatives form an algebraic basis of normal tensors for the original
vertical tensor by
\path{TNLean/MPS/MPDO/VerticalBNTConstruction.lean}. At every positive chain
length, this follows from exact reconstruction and cyclicity of trace. At
length zero, one retained representative of positive bond dimension supplies
the dimension identity required by the definition of a basis of normal
tensors.

Finally, \path{MPOTensor.verticalCF_of_horizontalCF} in
\path{TNLean/MPS/MPDO/VerticalCanonicalForm.lean} and the independent literal
theorem \path{MPOTensor.verticalCF_of_cpsvCanonicalForm} in
\path{TNLean/MPS/MPDO/CPSVVerticalCanonicalForm.lean} apply the shared
orthogonal-sector coisometry theorem. They identify the dependent sector pairs
$(\alpha,q)$ with standard finite direct-sum coordinates and prove
Eqs.~\ref{eq:cpgsv17_vertical_cf_grouping_source_coisometry},
\ref{eq:cpgsv17_vertical_cf_grouping_source_vertical_form},
and~\ref{eq:cpgsv17_vertical_cf_grouping_source_vertical_reconstruction}. The
source proposition does not print the reverse reconstruction separately. The
original formulation of \ghissue{2537}, interpreted as a direct conversion of
horizontal scalar weights and diagonal restrictions, remains a false route
rather than a consequence of Proposition~4.13.

\paragraph{Verdict.}
Proposition~4.13 of Ref.~\cite{Cirac2016MPDO_arXiv} is complete under its
literal CPSV canonical-form hypothesis. The theorem
\path{MPOTensor.verticalCF_of_cpsvCanonicalForm} produces a basis of normal
vertical representatives, positive diagonal multiplicities, and a rectangular
coisometry $U$ satisfying
Eq.~\ref{eq:cpgsv17_vertical_cf_grouping_source_coisometry}, together with
the forward and reverse identities
Eqs.~\ref{eq:cpgsv17_vertical_cf_grouping_source_vertical_form}
and~\ref{eq:cpgsv17_vertical_cf_grouping_source_vertical_reconstruction}. The
proof does not use the unsupported diagonal-restriction route from
\ghissue{2537}. The stronger theorem for
\path{MPOTensor.IsHorizontalCF} remains available separately. These
Proposition~4.13 results are distinct from the Appendix~C.4 positive-fusion
question tracked by \ghissue{4648} and from Theorem~4.14 of
Ref.~\cite{Cirac2016MPDO_arXiv}.

\section{Explicit BNT refinement and original-space Lemma L}

The preparatory dependency tracked by \ghissue{4956} is complete. The structure
\path{MPSTensor.CPSVCanonicalFormData.BNTRefinement} and its constructor
\path{MPSTensor.CPSVCanonicalFormData.exists_bntRefinement} retain an explicit
equivalence between phase-class copies and the originally listed blocks. Each
copy retains its original nonzero coefficient and bond dimension, together
with a unit phase, an invertible gauge, and an exact letterwise gauge-phase
relation to its representative. The representatives form a CPSV basis of
normal tensors.

The class-copy pairs form
\path{MPSTensor.CPSVCanonicalFormData.GroupedIndex}. The equivalence
\path{MPSTensor.CPSVCanonicalFormData.classCopyEquiv} identifies the grouped
order with the original block list. The equivalence
\path{MPSTensor.CPSVCanonicalFormData.groupedCoordinateEquiv} identifies the
flattened dependent coordinates with the original retained coordinates, and
\path{MPSTensor.CPSVCanonicalFormData.groupedTotalDim_eq} records the resulting
total-dimension equality. The reusable underlying results are
\path{MPSTensor.blockIndexCoordinateEquiv} and
\path{MPSTensor.toTensorFromBlocks_eq_reindex_blockDiagonal_equiv}.

For a grouped index $x$, let $\widehat\mu_x$ and $\widehat B_x$ be the weight
and block supplied by
\path{MPSTensor.CPSVCanonicalFormData.BNTRefinement.groupedWeight} and
\path{MPSTensor.CPSVCanonicalFormData.BNTRefinement.groupedBlocks}. The pair
$x=(\alpha,k)$ carries the original nonzero copy weight and the corresponding
phase multiple of $C_\alpha$. If $\Phi$ is the flattened coordinate
equivalence, define
\begin{align}
    T_{\mathrm{grp}}^i
        &=
        \operatorname{reindex}_{\Phi}
        \left(\bigoplus_x
        \widehat\mu_x\widehat B_x^i\right).
    \label{eq:cpgsv17_vertical_cf_grouping_completed_grouped_tensor}
\end{align}
This is the tensor
\path{MPSTensor.CPSVCanonicalFormData.BNTRefinement.groupedTensor}.

Let $G$ be the block-diagonal gauge on the originally listed blocks. The theorem
\path{MPSTensor.CPSVCanonicalFormData.BNTRefinement.groupedRegroupLetterwise}
proves, for every letter $i$,
\begin{align}
    \bigoplus_k\mu_kA_k^i
        &=GT_{\mathrm{grp}}^iG^{-1}.
    \label{eq:cpgsv17_vertical_cf_grouping_completed_regrouping}
\end{align}
The intermediate result
\path{MPSTensor.CPSVCanonicalFormData.BNTRefinement.regroupedTensor_eq_groupedTensor}
identifies the earlier in-place representative tensor with the explicit
grouped tensor. The theorem
\path{MPSTensor.CPSVCanonicalFormData.BNTRefinement.reconstructGrouped} then
uses the original ambient coisometry $U$, satisfying
Eq.~\ref{eq:cpgsv17_vertical_cf_grouping_source_coisometry}, to obtain
\begin{align}
    A^i
        &=U^\dagger GT_{\mathrm{grp}}^iG^{-1}U.
    \label{eq:cpgsv17_vertical_cf_grouping_completed_reconstruction}
\end{align}
Thus the class-copy coordinates, dependent-coordinate permutation, letterwise
block-gauge conjugation, and ambient reconstruction are explicit and exact.

The representative-sector decomposition and its positive-length
matrix-product-vector identity are explicit in
\path{TNLean/MPS/CanonicalForm/BNTRefinement.lean}. Its effective copy weight
is
\begin{align}
    \lambda_{\alpha,q}
        &=\nu_{k(\alpha,q)}\zeta_{k(\alpha,q)}.
    \label{eq:cpgsv17_vertical_cf_grouping_representative_effective_weight}
\end{align}
This weight is nonzero for every copy. The declaration
\path{MPSTensor.CPSVCanonicalFormData.BNTRefinement.representativeSectorDecomposition}
defines the representatives, copy counts, and weights. The theorem
\path{MPSTensor.CPSVCanonicalFormData.BNTRefinement.groupedTensor_sameMPV2Pos_representativeSectorDecomposition}
identifies the matrix product vectors of the full grouped tensor and the sector
tensor at every positive length. In this rendered identifier, the ordinary
digit~$2$ denotes the subscripted character used in the source declaration.

The simultaneous representative word-tuple span and its marked-chain
consequences are proved in
\path{TNLean/MPS/MPDO/CPSVRepresentativeGroupedLemmaL.lean}. In that file,
\path{MPSTensor.CPSVCanonicalFormData.BNTRefinement.eventuallyRepresentativeWordTupleSpan}
proves simultaneous representative word-tuple span at every sufficiently large
length.

The full-coordinate marked tensor uses the same listed coordinate space. Each
copy carries the marked block $\lambda_{\alpha,q}C_\alpha$. The exact trace
identity
\path{MPSTensor.CPSVCanonicalFormData.BNTRefinement.trace_groupedMarkedTensor_eq_representative_markedTensor}
states
\begin{align}
    \tr(\widehat C_{\mathrm{grp}}^sT_{\mathrm{grp}}^w)
        &=
        \sum_\alpha
        \left(\sum_q\lambda_{\alpha,q}^{|w|+1}\right)
        \tr(C_\alpha^sC_\alpha^w).
    \label{eq:cpgsv17_vertical_cf_grouping_representative_marked_trace}
\end{align}
The exponent $|w|+1$ counts the marked site. Eventual representative span then
proves the retained-coordinate marked and first-site forms of Lemma~L:
\path{MPSTensor.CPSVCanonicalFormData.BNTRefinement.groupedMarkedTensor_basis_eq_of_trace_agree}
and
\path{MPSTensor.CPSVCanonicalFormData.BNTRefinement.insertedTensor_representative_eq_of_firstSiteActionAgree}.
Neither conclusion requires an additional normalization hypothesis.

The original-space marked argument is complete in
\path{TNLean/MPS/MPDO/CPSVOriginalSpaceLemmaL.lean}. Suppose
\begin{align}
    A^i
        &=V^\dagger B^iV,
    \label{eq:cpgsv17_vertical_cf_grouping_original_space_reconstruction}\\
    VV^\dagger
        &=\Id,
    \label{eq:cpgsv17_vertical_cf_grouping_original_space_coisometry}\\
    C
        &=V^\dagger EV.
    \label{eq:cpgsv17_vertical_cf_grouping_original_space_mark}
\end{align}
Then, for every tail word $w$, including the empty word,
\begin{align}
    \tr(CA^w)
        &=\tr(EB^w).
    \label{eq:cpgsv17_vertical_cf_grouping_original_marked_coisometry}
\end{align}
The marked matrix makes the closed chain nonempty, so
Eq.~\ref{eq:cpgsv17_vertical_cf_grouping_original_marked_coisometry} does not
assert equality of unmarked length-zero matrix product vectors. The formal
result is
\path{MPSTensor.trace_marked_mul_evalWord_of_coisometry_reconstruction}; its
linear-marked specialization is
\path{MPSTensor.trace_linearMarkedTensor_mul_evalWord_of_coisometry_reconstruction}.

Linear marking of $T_{\mathrm{grp}}$ is exactly the full grouped marked tensor,
as proved by
\path{MPSTensor.CPSVCanonicalFormData.BNTRefinement.linearMarkedTensor_groupedTensor_eq_groupedMarkedTensor}.
Each copy carries $\lambda_{\alpha,q}$ times the representative mark.
Combining this identity,
Eq.~\ref{eq:cpgsv17_vertical_cf_grouping_original_marked_coisometry}, the
listed gauge, and the retained representative theorem gives literal-CF marked
Lemma~L on the original bond space:
\path{MPSTensor.CPSVCanonicalFormData.BNTRefinement.linearMarkedTensor_eq_of_trace_agree}
and
\path{MPSTensor.IsCPSVCanonicalForm.linearMarkedTensor_eq_of_trace_agree}.

For physical first-site actions, the same argument yields
\path{MPSTensor.CPSVCanonicalFormData.BNTRefinement.insertedTensor_eq_of_firstSiteActionAgree}
and its predicate-level form
\path{MPSTensor.IsCPSVCanonicalForm.insertedTensor_eq_of_firstSiteActionAgree}.
Positivity provides the required first-site agreement from a one-sided
invariant orthogonal projection through
\path{MPOTensor.firstSiteActionAgree_braRight_ketLeftBraRight_of_invariant}.
Consequently,
\path{MPSTensor.IsCPSVCanonicalForm.hasInvariantProjectorClosure_verticalTensor}
proves invariant-projector closure for the vertical tensor under literal CPSV
canonical form and \path{MPOTensor.IsMPDO}.

The literal pairwise Figure~8 identity is proved in
\path{TNLean/MPS/MPDO/CPSVFigureEight.lean} by
\path{MPSTensor.IsCPSVCanonicalForm.gramDressing_eq_of_two_grouped_corners}.
For a common normal representative,
\path{MPSTensor.IsCPSVCanonicalForm.gram_eq_pos_smul_gram_of_two_grouped_corners}
gives a real scalar $\omega>0$ such that
\begin{align}
    X^\dagger X
        &=\omega Y^\dagger Y.
    \label{eq:cpgsv17_vertical_cf_grouping_literal_pairwise_gram}
\end{align}
These theorems cover source lines~1909--1921 for any two corners satisfying
the source's positive similarity equations.

Literal Proposition~4.13 is therefore complete. The normalized physical-sector
maps are supplied by
\path{MPSTensor.IsCPSVCanonicalForm.exists_normalized_grouped_sector_maps}, and
\path{MPOTensor.verticalCF_of_cpsvCanonicalForm} constructs the final vertical
canonical form. Its rectangular map satisfies
Eq.~\ref{eq:cpgsv17_vertical_cf_grouping_source_coisometry}, not in general
$U^\dagger U=\Id$, and obeys both exact reconstruction identities
Eqs.~\ref{eq:cpgsv17_vertical_cf_grouping_source_vertical_form}
and~\ref{eq:cpgsv17_vertical_cf_grouping_source_vertical_reconstruction}. This
result is independent of the Appendix~C.4 positive-fusion question tracked by
\ghissue{4648}.
\bibliographystyle{alpha}
\bibliography{references}

\end{document}
